\section*{第 \ref{ch:b2:two-wave-interference} 章——双波干涉}

\begin{solution}{exo:b2:two-wave-interference:1}
$i = \lambda D/a = \qty{2.4}{mm}$；$x/i = 4.2$；第十条暗条纹位于 $9.5i = \qty{22.5}{mm}$ 处；
$a = \qty{2}{mm}$：\qty{0.47}{mm}；在水中为 $\lambda/n$：\qty{1.8}{mm}。
\end{solution}

\begin{solution}{exo:b2:two-wave-interference:2}
$1.25I_0 \pm I_0$：$V = 0.8$。$1.01I_0 \pm 0.2I_0$：$V = 0.2$——干涉项按
\emph{振幅}之比变化，为 $1/10$，而非强度之比。滤光后：
$V = 2\sqrt{0.5}/1.5 = 0.94$.
\end{solution}

\begin{solution}{exo:b2:two-wave-interference:3}
$2ne = \qty{665}{nm}$：亮纹位于 $665/(p + \tfrac12)$：\qty{443}{nm}（蓝紫色）；
暗纹位于 $665/p$：\qty{665}{nm}（红色）：薄膜呈蓝色。$e = \qty{1}{\micro m}$ 时：$2ne = \qty{2660}{nm}$，
亮纹位于 $760$、$591$、$484$、\qty{409}{nm} 处——四种颜色同时出现，
形成一片发白的淡粉色。
\end{solution}

\begin{solution}{exo:b2:two-wave-interference:4}
$r_p = \sqrt{p\lambda R}$：$1.09$、$1.53$、\qty{1.88}{mm}；$p = r^2/\lambda R = 85$，
即 \qty{1}{cm} 内的环数；中心为暗：厚度为零，发生一次符号反转。水中：
半径除以 $\sqrt n$，共 $98$ 个环。
\end{solution}

\begin{solution}{exo:b2:two-wave-interference:5}
(a) $i = 0.8$ 与 \qty{1.4}{mm}；第二条红色为 \qty{2.8}{mm}，第三条紫色为 \qty{2.4}{mm}。
(b) 第 $p$ 级红色位于 $1.4p$ mm 处，与第 $p + 1$ 级紫色位于 $0.8(p + 1)$ 处相遇：自
$p = 2$ 起各级发生重叠。(c) $\delta = ax/D = \qty{2.5}{\micro m}$：缺失的波长 $\lambda = \delta/(p +
\tfrac12)$：$714$、$556$、\qty{455}{nm}。(d) 中心为白色；随后的条纹带有
内侧蓝色、外侧红色的边缘；至第三级时，颜色
混合成白色。
\end{solution}

\begin{solution}{exo:b2:two-wave-interference:6}
(a) $b < \lambda L/4a = \qty{59}{\micro m}$。(b) $\pi ab/\lambda L = 6.7$：对比度为 $|\sin6.7|/6.7
= 0.06$。(c) 宽度为一厘米的光源会跨越数百个条纹间距的
位移，故不出现条纹。(d) $a < \lambda/\theta = \qty{61}{\micro m}$。
\end{solution}

\begin{solution}{exo:b2:two-wave-interference:7}
(a) $a = 2h = \qty{0.4}{mm}$：$i = \lambda D/a = \qty{1.25}{mm}$。(b) 在 $x = 0$ 处两路光程相等，
但该处条纹却为暗：反射附加 $\pi$（符号
反转，见于光密介质上的掠入射）。(c) $\delta = ax/D <
\qty{50}{\micro m}$：$100$ 个条纹。(d) 海面充当镜子，悬崖顶部的
天线充当屏：光源升高时，条纹掠过，
由此给出其仰角。
\end{solution}

\begin{solution}{exo:b2:two-wave-interference:8}
(a) $2d = 42\lambda$：$d = \qty{12.4}{\micro m}$。(b) $100/42 = \qty{2.4}{mm}$。(c) $2nd = p\lambda$：
$56$ 个条纹。(d) 条纹向远端移动，且随 $d$
减小而展宽；四分之一条纹相当于 $\lambda/8 = \qty{74}{nm}$ 的 $d$ 变化量。
\end{solution}

\begin{solution}{exo:b2:two-wave-interference:9}
(a) 仅第一次反射（空气 $\to$ 油，由疏到密）有附加相位；油 $\to$ 水的反射进入较低
折射率介质，无附加相位：$\delta = 2ne + \lambda/2$，亮纹位于 $2ne = (p + \tfrac12)\lambda$。(b) $2ne =
\qty{1160}{nm}$：增强的为 $773$（可见光边缘）、\qty{464}{nm}（蓝色）；
抵消的为 $1160/p$：\qty{580}{nm}（黄色）：呈蓝色。(c) $2ne = \lambda/2$：$e = \qty{72}{nm}$。
(d) $\delta = 2ne\cos r$ 随角度减小，增强的波长
向蓝端移动。(e) 此时两次反射均使符号反转，不再附加
$\lambda/2$：亮纹位于 $2ne = p\lambda$，暗纹位于 $(p + \tfrac12)\lambda$——即
四分之一波长增透膜。
\end{solution}

\begin{solution}{exo:b2:two-wave-interference:10}
(a) $(n - 1)e = \qty{10}{\micro m}$，$20$ 个条纹，向被覆盖的狭缝一侧移动（此处
几何光程须缩短 \qty{10}{\micro m}，以恢复 $\delta = 0$）：
$x = (n - 1)eD/a = \qty{2}{cm}$。(b) 白色条纹位于该处；玻璃的
色散（可见光范围内 $\Delta n \approx 0.01$，即光程 \qty{0.2}{\micro m}）
小于 $\ell_c$：仍可见，略带颜色。(c) 测量白色条纹的
位移：$e = (\text{位移})\,a/(n - 1)D$。(d) $e(n - 1)\theta^2
(1 + 1/n)/2 = \qty{63}{nm}$，在 \qty{5}{\degree} 时：相当于八分之一条纹。
\end{solution}

\begin{solution}{exo:b2:two-wave-interference:11}
(a) $\theta = \qty{2.3e-7}{rad}$：$a = 1.22\lambda/\theta = \qty{3.1}{m}$。(b) 其极限 $1.22\lambda/d =
\qty{0.057}{''}$ 大于恒星的角直径，大气则将其模糊至
$1''$。(c) $1.22\lambda/100 = \qty{7e-9}{rad} = \qty{1.4}{mas}$。(d) 条纹需满足
$\delta < \ell_c = \lambda^2/\Delta\lambda = \qty{32}{\micro m}$.
\end{solution}

\begin{solution}{exo:b2:two-wave-interference:12}
(a) $\delta = 2ne\cos r + \lambda/2$；$\cos r = \sqrt{1 - \sin^2i/n^2} \approx 1 - i^2/2n^2$。
(b) $2ne/\lambda = 10186.8$：中心处，$p + \tfrac12$ 介于 $10186.5$ 与
$10187.5$ 之间，既非亮纹也非暗纹。(c) $i_p^2 = 2n^2[1 - (p + \tfrac12)\lambda/
2ne]$，对 $p = 10186$、$10185$、$10184$，有 $i = 0.011$、$0.024$、\qty{0.032}{rad}，$r = fi
= 5.4$、$11.8$、\qty{15.8}{mm}。(d) 仅角度 $i$ 即可确定级次：每个
光源点都以相同角度向同一圆环供光，因此这些圆环在无穷远处观察时，
即使光源较宽仍然保持。
\end{solution}

\begin{solution}{pb:b2:two-wave-interference:1}
\textbf{1.} 由对称性，每面镜子使 $S$ 成像于距接触线 $d$
处；两个像位于半径为 $d$ 的圆上，
相隔角度 $2\alpha$：$a = 2\alpha d = \qty{2}{mm}$。

\textbf{2.} $S_1$ 与 $S_2$ 是同一振动的两个像：其
相位跃变相同。

\textbf{3.} $i = \lambda(d + D)/a = 589 \times 10^{-9} \times 1.7/2 \times 10^{-3} = \qty{0.50}{mm}$：
在 \qty{5}{mm} 内有十个条纹。

\textbf{4.} $\delta = ax/(d + D) = \qty{2.9}{\micro m}$（边缘处），$p = 5$：远小于
$\ell_c/\lambda = 34\,000$.

\textbf{5.} 位移 $bD'/d = 20 \times 10^{-6} \times 1.7/0.2 = \qty{0.17}{mm}$，为
$i$ 的三分之一：对比度为 $|\operatorname{sinc}(\pi ab/\lambda d)| = \operatorname{sinc}(1.07) = 0.82$。

\textbf{6.} $\pi ab/\lambda d = 5.3$: $|\sin5.3|/5.3 = 0.16$.

\textbf{7.} 振幅为 $\sqrt{0.9}$、$\sqrt{0.8}$：$V = 2\sqrt{0.72}/1.7 = 0.998$。

\textbf{8.} $(n - 1)e = \qty{5}{\micro m}$：$8.5$ 个条纹，钠光下极易跟踪；
白光下白色条纹移动 $8.5i = \qty{4.2}{mm}$，
仍清晰可见（玻璃的色散只造成
不足一微米的偏移）。

\textbf{9.} 中央为白色条纹，两侧各有两三条彩色条纹，
其外为均匀的白光。

\textbf{10.} 空气 $\to$ 油时发生反转，油 $\to$ 水（折射率较低）时则不然：
$\delta = 2ne + \lambda/2$；亮纹条件为 $2ne = (p + \tfrac12)\lambda$。

\textbf{11.} $e = \lambda/4n$：红色 \qty{112}{nm}，绿色 \qty{95}{nm}，紫色 \qty{72}{nm}。

\textbf{12.} $2ne = \qty{1740}{nm}$：增强的为 $1740/(p + \tfrac12)$：$696$（红色）、$497$
（蓝绿色）、\qty{387}{nm}；相消的为 $1740/p$：$580$（黄色）、\qty{435}{nm}：
呈紫红色，是红色与蓝绿色的混合。

\textbf{13.} 膜厚沿路径变化；每种颜色在 $2ne$ 取适当值时
为亮，因此这些条带是
等厚线。

\textbf{14.} 级次递减：颜色依次沿序列
向薄膜端变化；当 $e \to 0$ 时，唯一的一次反转使
各色均有 $\delta = \lambda/2$：暗纹——与空气中肥皂膜的情形一样。

\textbf{15.} 厚度仅几纳米：各色均有 $2ne \ll \lambda$，两束
反射光反相而相消：呈黑色。

\textbf{16.} $\sin r = \sin60^\circ/1.45$, $r = \qty{37}{\degree}$, $\cos r = 0.80$: $2ne\cos r
= \qty{1390}{nm}$，亮纹位于 $928$、$557$、\qty{398}{nm} 处：第 12 题中的红色已
变为黄绿色——趋向蓝色。

\textbf{17.} 当 $2ne$ 超过 $\ell_c \approx \qty{1}{\micro m}$，即油层厚度超过
约 \qty{0.4}{\micro m} 时，各级次重叠，颜色褪为灰色。

\textbf{18.} 对任一光源方向，由膜面反射的两束光
均在膜面上相遇：条纹定域于膜面，
宽而多方向的光源仅是叠加相同的图样。

\textbf{19.} $\ell_c/\lambda$：用钠光时为 $34\,000$，用激光时为 $1.7 \times 10^7$。

\textbf{20.} $b < \lambda d/4a$：平面镜情形为 \qty{15}{\micro m}，
狭缝情形为 \qty{59}{\micro m}。

\textbf{21.} 画出两个光源点：对每个光源点，两束反射光
在膜面的同一点会合，并具有相同的 $\delta = 2ne\cos r$
（近似垂直入射）；膜的图样不动，杨氏条纹却
会移动。

\textbf{22.} $\lambda/\Delta\lambda = 589/0.02 \approx 30\,000$.

\textbf{23.} $550/30 \approx 18$ 个条纹。

\textbf{24.} 相干光经屏幕颗粒散射，在眼睛处以随机相位
发生干涉：形成散斑。

\textbf{25.} 平面镜：分波前，条纹处处可见
（非定域），条纹数受 $\ell_c/\lambda$ 限制，宽光源会将其
抹去。膜：分振幅，条纹位于膜面上，颜色
数受 $\ell_c \sim \qty{1}{\micro m}$ 限制，宽光源无妨。
\end{solution}
